\section{Discussion}\label{sec:discussion}

\subsection{Interpretation}\label{sec:interpretation}

The price and volume results point in similar directions but identify different things. On aggregate, prices for agreement countries exhibit higher initial pass-through followed by partial absorption, suggesting a moderating effect of trade agreements on tariff transmission to border prices. This pattern holds across both steel and non-steel subsamples, but pre-trend violations within both subsamples preclude causal interpretation of the price differential, potentially driven by fundamental differences in the composition of products imported from agreement versus non-agreement countries. Volumes, by contrast, satisfy parallel trends and admit a causal reading. The aggregate volume results reveal a large positive differential agreement effect, driven by rising import volumes from agreement countries and falling volumes from non-agreement countries. However, this aggregate finding masks meaningful heterogeneity: the volume gain for agreement countries is driven by the steel subsample, while non-steel agreement countries experience only a transient positive effect that converges back to non-agreement levels after roughly six months. Analyzing the steel and non-steel subsamples revealed pre-trend violations that were not statistically detectable in the aggregate sample. While these findings are consistent with the hypothesis that trade agreements are resilience-enhancing in the face of unilateral tariff shocks, the difference in identification quality between prices and volumes, as well as the product-category heterogeneity within volume elasticities, mean that causal interpretation must be at the product-category level and rely primarily on volume evidence. 

One possible mechanism behind the volume phenomenon could be that pre-existing trade agreements provide a buffer in the form of prior tariff and non-tariff barrier reductions. These leave agreement-country exporters with greater margin to absorb new tariff costs, resulting in lower pass-through of tariffs to prices, making agreement-country firms more competitive and thus reallocating demand towards them. This is consistent with the aggregate finding that agreement countries experienced slightly smaller pass-through, though the price results were insufficient for clear identification. 

Two expectation-based channels may complement this effect: greater expected stability due to agreements and anticipated exemptions. Importers facing an uncertain trade policy environment might rebalance along the intensive margin, increasing imports from agreement-countries that have more certain environments while shifting away from non-agreement-countries, consistent with relative rebalancing identified in \textcite{carballo2022}. Meanwhile, as the Section 232 investigations that justified tariffs were publicized well in advance of tariffs' implementation, importers may have made arrangements to import more from markets that were more likely to receive exemptions, such as Canada and Mexico, with which the US has a bilateral trade agreement. The eventual reconvergence after six to seven months may reflect either renewed relationships with previous suppliers from non-agreement markets, or the substitution of similar suppliers from the same non-agreement markets. 

Product-specific characteristics may affect how particular products respond to tariffs. Steel products' homogeneity may mean that importers can easily substitute steel products sourced from agreement countries, locking in those suppliers for extended periods of time without significant concern, as stability is more important than price or product characteristics. Non-steel products, which in this sample contain solar panels, washing machines, and aluminum products, may not be so easily substitutable. Two effects may influence non-steel products' transient effect. First, there may be fundamental differences in non-steel products from agreement versus non-agreement countries that make those products from agreement countries unfavorable in the longer-term, such as quality or differentiating factors. Second, it may be the case that agreement countries simply do not have sufficient capacity to satisfy US demand for such products, and therefore can only serve as substitutes in the short-term. Simultaneously, aluminum is likely to behave similarly to steel, as it is also a homogenous input. Given it is likely a large component of the non-steel sample, there may be some attenuation bias in the differential effect of tariffs on non-steel volumes between agreement and non-agreement countries. 

While the evidence points to agreement countries substituting non-agreement countries, the data does not describe the magnitude of this effect. It is clear that, particularly for steel products, volume rises following tariffs for agreement countries while it declines for non-agreement countries. However, this data does not reflect the starting points of these products' import values, only the log changes, interpretable as percent changes. Particularly because the agreement sample has substantially fewer country-product observations, it is likely that volumes are also substantially smaller. Though the elasticity of steel import volumes for agreement countries rises by as much as --- if not more --- than the non-agreement countries' imports fall, the level of volumes is unlikely to be completely offset. Therefore, this implies that overall import volumes fall as a result of tariffs, as supported by the aggregate findings in \textcite{amiti_whos_2020}. 

The differential effect of trade agreements on price pass-through due to tariffs is ultimately not identifying, but is useful. While differential price elasticities must be interpreted as descriptive rather than causal due to identification quality issues, they are useful in order to contextualize volume elasticities. The trend of partial tariff absorption within agreement-country products that is indicated by the aggregate differential price elasticities supports the divergent pattern seen between agreement and non-agreement countries in volume elasticities. 

\subsection{Limitations}

This study uses a triple-difference design, finding the difference between the pass-through coefficients for countries with and without trade agreements with the US over time. This approach might confound the results, as the fixed effects are applied using a different sample set. As a result, the results could have some form of bias should the samples meaningfully deviate in some dimension. 

USMCA was signed in 2018 and entered into force in 2020. While this means it was not in force during the period of this study, anticipatory effects based on policy differences between USMCA and NAFTA may have influenced results. 

\subsection{Potential robustness checks and extensions}

There are several robustness checks worth pursuing. First, trade literature suggests that a Poisson Pseudo-Maximum Likelihood (PPML) estimator may produce more accurate results by accommodating for zero-trade values, which the OLS estimator requires us to remove. This may attenuate treatment effects. Second, since Canada and Mexico received tariff exemptions in May 2019, implementing the model with and without including them would test the results. Third, given the triple-difference interpretation, a well-specified triple-interaction term would test and confirm results' accuracy, allowing for a more accurate comparison across groups, as fixed effects would be defined on a common sample. Fourth, given Jordan is the one country without a trade agreement considered deep by my approach, adding or removing it from the sample, or testing it individually would permit some examination of whether depth does play a factor. Fifth, pooling data from months into quarters ...

There are several possible extensions to this study that would further illuminate the differential effect of trade agreements on the impacts of tariffs. First, further heterogeneities (e.g. across waves, end-use categories) may reveal interesting patterns. Second, theory suggests small and large firms respond differently to international economic shocks --- further research could investigate whether the distribution of trade flows by size differs across agreement and non-agreement countries as a result of tariffs.  